
% Chapter 1 File

\chapter{Introduction}
\label{introduction}
\thispagestyle{empty}


The mobile industry has incredibly changed throughout the years it existed. While the developments in hardware and software components have been in the public eye for many years now, the most impressive change was arguably the shift in power and influence among market participants. 
In fact, the industry's balance of power, which back in the 1990s and early 2000s used to favour mobile device manufacturers and, in particular, network operators, now sees a hardly disputed winner, which are mobile platform providers, like Google or Apple (Basole et. Karla, 2012, p. 35).

The most important year in this transition was 2008 and Apple was probably the main actor. As commonly known, in 2007, Apple launched the first version of the iPhone and it could be argued that the real innovation didn't lay in the new hardware features of this phone: Apple's iPhone wasn't the first one to introduce a capacitive touchscreen, for example, and LG did that some time in advance in 2006 (Ihlwan, 2007). The real breakthrough that initiated the ``app economy" was the introduction, in 2008, of Apple's App Store, a mobile application marketplace with originally a few more than 500 apps and which now accounts for more than 1.5 million applications (Basole and Karla, 2012, p. 32; Ceci, 2022b; Friedman, 2013). The first app stores were launched at the end of the 1990s, but it wasn't until almost a decade later, with the Apple Store, that they started to gain traction (Kouris and Piller, 2014, p. 21). Since then, the increasing availability of mobile apps, supported by enhancements in hardware features, drove the growth of the mobile industry. In fact, the number of smartphone users worldwide almost duplicated only considering the last five-year period, reaching 6.2 billion users in 2021 and being expected to grow up to 7.6 billion in 2027 (Ericsson, 2022). Meanwhile, apps became an incredibly profitable business by themselves and those downloaded from the Google Play Store, the largest app store by the number of available apps at the moment (3.5 million apps in 2022), generated in 2021 more than \$ 47 billion in gross revenue, three times more than what was generated back in 2016 (Ceci, 2022a, 2022b). 

Clearly, being the owner of a successful app marketplace is a position as desirable as difficult to reach for companies operating in the mobile ecosystem and many appear to fail while attempting that. The complex dynamics lying behind the interaction among market participants (e.g. app developers and customers) are shaped by the mechanics of information economics, which may raise entry barriers and uncertainty. While past literature has already provided useful insights on two-sided platforms,\footnote{Since the word ``platform" may refer to different elements of the software ecosystem, ``app (or software) store or marketplace" will be used from now on referring to platforms from which software can be purchased and downloaded, as, for example, Google Play Store or Apple App Store for apps or Steam for PC videogames. On the other side, ``technology platform" or just ``platform" will refer in a broader sense to all the software and hardware elements that allow an application to run on a device, often including as its fundamental part the operative system (OS). Finally, ``two-sided platform" may be used to generically refer to any two-sided market, beyond those dealing with software applications.} a category to which app marketplaces belong, some elements of customer behaviour, when dealing with digital products, need to be explored more deeply. However, before investigating users' preferences on mobile app stores and focusing on empirical evidence, an overview of the mobile ecosystem is necessary.


\section{Mobile Ecosystem}

The mobile phone industry is now more than forty years old. In 1978, in fact, a Motorola engineer, Martin Cooper, who is today considered the inventor of the cell phone, made the first mobile phone call using a Motorola DynaTac in New York (Zheng and Ni, 2010, p. 32). Six years later, in 1984, a more recent version of the same Motorola phone was commercialized for the first time. Those phones were also defined as ``brick phones", due to their shape and size, and, in the following ten or twenty years, the innovation efforts of this industry focused on miniaturization, better power efficiency and battery life, better displays and signal. In the meantime, generations of cellular network technologies were succeeding (from 1G to today's 5G), improving voice communication quality, as well as introducing revolutionary features such as text messaging or Internet access. However, the real disruption in this industry came in 2008, with the introduction of the first smartphones and of the app concept, which overturned power relations between market participants. 

Basole and Karla (2012, p. 28) identified four groups of market participants, mobile platform providers, mobile device manufacturers, mobile network operators, and mobile application providers, to which a fifth group, mobile users, can be added to effectively understand the industry's internal dynamics.

\paragraph{Mobile Platform Providers.}

Mobile platform providers (MPPs) provide operating systems and middleware solutions for smartphones and they are currently driving the mobile industry. In addition, MPPs often coincide with the owners of major mobile app stores. This double role was arguably assumed for the first time by Apple in 2008, when, while launching the iPhone, it leveraged its experience in content distribution, thanks to the previous existence of iTunes, to introduce the App Store, in order to strengthen its ecosystem and drive demand for iOS devices (Basole and Karla, 2012, p. 36; Cortimiglia et al., 2011, p. 3), before app stores became a profitable business by their own.  While a couple of years after its introduction, iOS was installed on 37 percent of all smartphones, a market share similar to the one of Symbian OS (Nokia's own operative system), and Android had just entered the market, as shown in Fig. 1.1, now the latter dominates other mobile operative systems (with about 71 percent of the global market), just followed by iOS and other operative systems with very small shares (Laricchia, 2022).
\begin{figure}
		\includegraphics[width=\linewidth]{Figure/Figure 1.png}
		\caption{Mobile OS market shares in 2022 and 2009 (Laricchia, 2022).}
		\label{fig:graph1}
\end{figure}Interestingly, these two systems were both able to succeed with two different approaches: iOS is considered a relatively closed ecosystem (``walled garden"), while Android is a more open one (``open garden"). The unique role of Apple as both MPP and device manufacturer enables it to apply some restrictions, as, for example, the fact that apps can only be installed from the App Store, as sole source (Basole and Karla, 2012, p.36; Cortimiglia et al., 2011, p. 3), and this approach allows more strict testing and verification, that, in combination with the existence of only a few different hardware with that operative system (OS), may lead to a better user experience, with consequently improved customer loyalty and a premium brand image. On the other hand, Android is an open-source OS and the ecosystem is more open: this may lead to higher levels of competition both on software and hardware, with lower prices as result (Müller et al., 2011, p. 74). In terms of app distribution, both Apple and Google act as gatekeepers: they decide what gets published on their stores, also covering important roles in terms of quality assurance and pre-screening of dangerous apps (Müller et al., 2011, p. 73), but, while in Google case, these limits can be circumvented thanks to the presence of other app sources, this is (legally, at least) impossible on Apple's devices.

Overall, MPPs took over many parts of the value chain, from OS development to payment systems, relegating other participants to more subordinate roles.


\paragraph{Mobile Device Manufacturers.}
Mobile device manufacturers' (MDMs) activity mainly consists in producing and distributing the hardware part of a smartphone. In the past, MDMs generally owned proprietary operative systems installed on their products, only a few native apps were offered with the phone and the possibility of installing third-party apps was generally limited. Their role was strictly connected to the one of network operators, on which they often depended for distributing devices, usually sold in combination with mobile subscriptions (Basole and Karla, 2012, p. 29). 

Today, the only MDMs developing a popular OS are the most important MPPs, Apple and Google, while all the other producers (Samsung, Xiaomi, Oppo, Vivo, etc.) redirected themselves towards Android as the chosen OS for their devices in most cases.


\paragraph{Mobile Network Operators.}
Mobile network operators (MNOs) have been dominant in the mobile service industry for many years. For decades, they were the main sources of innovation and added value, they supplied SIM cards and they could also dictate some pre-settings, like the browser's homepage or types of applications installed (Basole and Karla, 2012, p. 29). Later, with the smartphone introduction and the app economy emerging, MNOs lost their traditionally central role, their margins shrunk and the risk of becoming dumb pipes is real.


\paragraph{Mobile Application Providers.}
Mobile application providers (MAPs) are those providing content in the current mobile ecosystem. They develop both web-based and native applications for specific technology platforms, which cover many different categories and use cases. Developers traditionally lacked market accessibility and app stores greatly reduced barriers and time-to-market for them, but, on the other side, they now largely depend on MPPs' enabling capabilities and developing tools (Basole and Karla, 2012, p. 30). Anyway, MAPs can now reach a large number of potential customers, at the price of leaving (on average) 30 percent of their app download revenue to gatekeepers (Cortimiglia et al., 2011, p. 5), whose role is particularly restricting in walled gardens, due to the lack of alternative channels.


\paragraph{Mobile Users.}
On the other side of a mobile app store, there are smartphone users. They search, download and pay for apps and they also later rate them. There is a tendency among users to view the mobile device and available apps as a package: they take into consideration the ecosystem as a whole and then decide which hardware and software combination to buy. In this sense, mobile devices can be considered complementary goods for apps and their price a one-time membership fee to join the ecosystem (Kouris and Piller, 2014, p. 36).

For customers, downloading an app and the perception of the risk involved depend more on the MPPs than MAPs: users, in fact, place their trust in platform owners as their main source of information. In addition, they use apps to perform highly fragmented everyday tasks and perceive apps not as complex as traditional desktop software. This makes their buying decision limited, habitual and impulsive, corresponding to generally low levels of involvement (Buck et al., 2014, p. 5). During the download process, users often don't search for information actively and reputation systems (e.g. rating score), as well as peer influence (e.g. sale rankings, number of downloads), drive their decision-making (Buck et al., 2014, p. 5; Liu et al., 2014, p. 330). 


\section{Mobile Apps as Digital Products}

The term ``mobile application", ``mobile app" or just ``app" refers to application software, which is ``a type of software that allows the user to perform a specific task, that can be installed and run on a range of portable devices such as smartphones and tablets" (Liu et al., 2014, p. 327). An app usually distinguishes itself by general software because it is optimized for touch-screen-based mobile devices, it is offered free of charge or for sale and the most of the times it is available for download through a centralized online marketplace where users can also rate and review the app and access some ranking lists (Liu et al., 2014, p. 327). Apps already have a longer history than one could expect: for example, Zheng and Ni (2010, pp. 51-56) were able to individuate two generations of mobile apps already back in 2010. The first generation involved very simple apps, such as ringtones, instant messaging and personal information management apps, among which the first one was the most popular for many years. Then, the second generation included more complex apps like mobile commerce, gaming, music streaming and mobile social networking and some of these are still popular today. Later, a third generation emerged, with features like augmented reality, ubiquitous social networks and merging and seamless switching between different functions (Kouris and Piller, 2014, p. 25). 


Mobile apps, exactly like all types of software, are digital goods and as such they share all the typical characteristics of this category of products. Compared to standard economic models, in a digital information setting, costs fall substantially and approach zero. In fact, Goldfarb and Tucker (2019, p. 3) were able to identify five types of cost reduction: lower search costs, lower replication costs, lower transportation costs, lower tracking costs and lower verification costs. Even if these cost reductions are intrinsic to all digital (i.e. intangible) goods and services, they are enhanced by the use of mobile app stores as the main distribution channel. 

In fact, mobile app stores are an example of two-sided platforms. These platforms are intermediaries that enable exchanges between other players and, by doing that, they efficiently organize supply and demand allowing buyers to access the sellers (and vice versa), facilitating matchmaking and lowering transaction costs. Platforms aggregate information and provide access to that data to market participants and also provide ratings and feedbacks that help to build credibility and trust among participants, making some previously unknown seller characteristics observable, rewarding the better-quality users and creating incentives for lower-quality ones to leave the market. In exchange for easier transactions, these businesses charge a share of their value for the service (Lambrecht et al., 2014, p. 10). In addition, by aggregating market participants, intermediary platforms generate indirect externalities (or indirect network effects): in fact, they usually consist of two-sided (or multi-sided) networks with two (or multiple) types of members, where each type is looking for a network with the largest number of users of the other type. In these cases, network effects become so strong that customers tend to pick a network, instead of choosing a service or a product, since the utility on one side of the market highly depends on the number and quality of participants on the other side (Kouris and Piller, 2014, pp. 39-40; Piccoli and Pigni, 2019, p. 239). Indeed, mobile app stores are characterized by a positive feedback loop: the greater the number of users of the store (reflected in the number of downloaded apps or in the installed subscriber base), the more attractive that marketplace becomes for developers, resulting in an increase of the supply of content and applications and, as a consequence, in higher consumer benefits and increased attractiveness for additional users (Basole and Karla, 2012, p. 35).

\paragraph{Lower Search Costs.} 
It is a lot easier to find and compare information about potential transactions online than it is offline. As a consequence of the decrease in search costs, prices have fallen and product variety increased (Goldfarb and Tucker, 2019, pp. 6-7). While prices online appear to be constantly lower than offline, also for tangible products (e.g. Brynjolfsson and Smith, 2000, p. 579), in purely digital contexts, such as app marketplaces, prices actually reached zero in many cases. Minimal search costs, together with other reasons that will be later discussed, had an important role in spreading and making successful the freemium strategy and increasing the offer of free apps.

However, even if prices are lower, an important level of price dispersion persists online. The ease of changing prices and running price promotions on mobile apps, which also have no marginal costs of production, makes this a common characteristic of mobile app stores too. ``Free App of the Day" promotions are an example of an initiative that could rise price dispersion on app marketplaces. This kind of initiative involves spotlighting in the homepage and zero-pricing an app for 24 hours and it had been used in the past, for some periods, by the Google Play Store, the Apple App Store and the Amazon Appstore, the three most popular app stores today (Chen et al., 2022, pp. 1-2). The objective of this practice is combining digital free sampling, which proved to be effective in reaching new customers, with sponsorship practices, like featuring the app in a noticeable position, consequently increasing visibility and enhancing customer awareness of the app. According to Chen et al. (2022, pp. 6-7), such a promotion may lead to adverse effects in the post-promotion phase: firstly, the removal of zero-price hurts sales, in particular, due to customers' expectations of the app being free and their perception of a loss by observing the restored (and increased) price; secondly, promotion-induced reviews, generated by the increased visibility of the app, are less effective than organic ones.

Low search costs also mean that it is way easier to find rare and niche products online (Goldfarb and Tucker, 2019, p. 8), a phenomenon defined by Anderson (2006, p. 22) as ``the long tail". The effects of the long tail are particularly visible on app stores: apps generally require just relatively little storage space, making welcoming as many apps as possible quite inexpensive for marketplace owners, with positive benefits overall. Regardless of the number of sales, unsuccessful apps provide additional choices and options for customers, reinforcing the two-sided positive feedback loop (Basole and Karla, 2012, p. 35). In addition, the growing user-friendliness of developing tools allowed many people besides professional developers to publish their apps, further increasing the range of available options. A broader choice of apps, from a negative perspective, leads to higher choice complexity, which could cause delays in the users' app download process. This negative effect of the long tail is usually mitigated by recommendation systems, app rankings and improved research, made possible by the available data on customers' previous behaviour.

Finally, an online marketplace can sometimes generate the superstar effect, with homogeneous customers all agreeing on which product is the best and buying it (Goldfarb and Tucker, 2019, p.8). In the case of apps, this phenomenon is extremely strengthened when software benefits from direct network effects and its utility largely depends on the size of the installed base (e.g. instant messaging apps). In these cases, the probability of being in presence of a tipping (or ``winner-takes-all") market is high.

\paragraph{Lower Replication Costs.}
In addition to being subject to marginal costs which are close to zero, software is a non-rival, meaning that it can be consumed by one person without reducing the amount of quality available to others, and, in absence of a legal or technological limit, non-excludable good, so it can be reproduced and used by anyone at approximately no cost. These characteristics often generate opportunities, like the possibility of implementing bundling strategies at a very low marginal cost, very common in the software industry, but also some threats.

Low replication costs and non-excludability are piracy's main reasons for existence. Piracy has been a long-discussed issue in the software industry and illegally distributed copies actually offered in some ways opportunities for free sampling and strengthening network effects in the past, leading researchers and companies to think of piracy as an advantage under certain conditions. Raising the barriers of protection causes, in fact, some users to buy the product and others to live without it or look for alternatives, consequently reducing the installed base and its value also for those users that paid for the program (Reavis Conner and  Rumelt, 1991, pp. 136-137). In addition, illegal copies contribute to raising switching costs for pirates, with a higher probability that they would later buy the software. Today, the incentive not to get in the way of illegal distribution has been partially revised, also considering the effectiveness of freemium software in covering the same functions previously fulfilled by piracy, while remaining in control of intellectual property and generating some earnings from advertising, for example (Halmenschlager and Waelbroeck, 2014, pp. 30-31; Nan et al., 2018, p. 283).

In some way, compatibility can be seen as a limit to software duplication: in fact, software needs to be compatible with other elements of the ecosystem (the operative system, in particular) and, therefore, in some cases, it cannot be run without some efforts by developers to make this possible. Sometimes, making software compatible with an additional OS can be almost as expensive as it was creating it at first and it needs some incentives from the platform owner side, such as, for example, a large installed base, a high-quality development toolkit or low monetary and technical entry barriers: developers that choose to target two or more ecosystems try to balance between a greater potential market and the cost of porting the product for different platforms (Hyrynsalmi et al., 2012, p. 60; Koch and Kerschbaum, 2014, pp. 1432-1434; Rochet and Tirole, 2003, pp. 1015-1017). Examining the strategy of the most popular apps, multi-homing might be assumed sometimes to be a viable way to increase the circulation of an app. However, the effect could be bi-directional and popularity in a single ecosystem could actually trigger the porting of the app in a second one. In fact, on mobile app stores, the large majority of apps (about 87 to 95 percent) are developed for a single technology platform and apparently, a multi-homing strategy isn't very effective in making an app more popular than others (Ali et al., 2017, p. 5; Hyrynsalmi et al., 2012, p. 71). While at least professional developers would be expected to develop cross-platform apps, one out of five top-rated apps are not present on both Android and iOS, for reasons which reportedly are lack of resources, platform restrictions, low expected revenue in the other platform, user fragmentation within a platform and competition being already present, and, even if developed for both platforms, in any case, they are usually not consistently released in terms of new versions and updates (Ali et al., 2017, pp. 6-7). Obviously, if multi-homing can be hard for companies with skilled developers, it gets even rarer in an ecosystem that contains a good number of app publishers doing it just for fun or lacking a clear monetisation plan (Hyrynsalmi et al., 2012, p. 68). In addition to that, platforms in some industries tend not to promote multi-homing and even to obstacle it when it reduces possibilities of differentiation. While setting conditions that raise developer attractiveness usually increases the share of cross-platform software available, it also reduces user attractiveness when compared to other platforms (Hyrynsalmi et al., 2012, p. 70).

\paragraph{Lower Transportation Costs.}
The cost of transporting information goods over the Internet is quite near zero. Software has been distributed online and in physical formats for many years now and, while, back in the days, offline distribution was certainly more widespread, the Internet started to pick up as an important distribution channel soon.

Anyway, software distribution over the Internet used to present some drawbacks, in particular in the past. In fact, anyone could post software to the Internet, with no accountability and making malicious software downloads quite common, creating the need for a solution for trusted software distribution (Rubin, 1995, p.7). From a certain perspective, app stores answered to that problem, creating a trusted and secure channel even for less technologically advanced customers. In addition, they also provided a more efficient and faster system for automated software updates, making them almost unnoticeable to unaware customers, reducing security threats and improving user experience.

Today, when getting information on which app to download, users place their trust in app store owners, creating a way closer relationship with intermediaries than third-party software developers (Buck et al., 2014, p. 13).


\paragraph{Lower Tracking Costs.}
In a digital domain, user activity can be easily recorded and stored (Goldfarb and Tucker, 2019, p. 19), enabling relevant opportunities to make software selling more efficient. A major example of this is personalized recommendations, present on almost every online marketplace as well as on mobile app stores. 

The availability of information on user behaviour also opened up new price discrimination opportunities: in particular, freemium, a specific kind of versioning, became predominant. Freemium, a combination of the words ``free" and ``premium", is a business model by which a service or a product is offered free of charge, but a premium version is charged for advanced features, functionality, related products or services (Liu et al., 2014, pp. 327-328). Therefore, the free basic version of the product contains limited features, less privacy, more advertising or limited time usage (Boudreau et al., 2022, p. 1378). Freemium insert itself in a context of sampling and versioning strategies, that proved, for many years, to be efficient in  creating a positive brand attitude and awareness when uncertainty about the product is high (Arora et al., 2017, p. 63). In addition to that, mobile applications are largely experience goods and it can be difficult for customers to quantify the true value of the product before using it, while a free version of the software allows customers to get to know the product, expecting that a considerable share will purchase the premium version (Boudreau et al., 2022, p. 1379). The freemium strategy in the app economy is also used to exploit advantages ascribed to positive network effects, product complementarity and an accelerated diffusion process. For companies leading the market, a free version of an app may provide a method for exploiting their advantage and expanding their lead over their competitors; for new entrants, a freemium strategy may provide a way to eat a part of the leader's advantage (Boudreau et al., 2022, pp. 1375-1376). However, unlike traditional software, apps are generally characterized by minimal learning and switching costs and a low degree of interdependence among apps from the same developer (Liu et al., 2014, p. 330), raising fragmentation and competition. Implementing this strategy remains anyway quite complex and a trade-off emerges: on one side, consumers will judge the app on the basis of the quality of the free version (Liu et al., 2014, p. 330); on the other side, if the free app is sufficient for users, it could cannibalize the paid version.

Nevertheless, collecting users' information isn't always positive and privacy regulation puts a cost on tracking information flows (Goldfarb and Tucker, 2019, p. 23): on one hand, there are legal compliance costs and, on the other, there is a huge reputation risk. Many users report being very concerned about their privacy, but they still frequently download free apps that monetize their personal data, bringing to light a piracy paradox. Clearly, consumers on the Internet articulate their need for privacy, but they act in an opposite direction (Buck et al., 2014, p. 4).
 
\paragraph{Lower Verification Costs.}
As already discussed, mobile app stores, as two-sided platforms, contribute to increasing trust among participants through the rating system. In an offline setting, the company's reputation rests on having an uninterrupted series of transactions with a positive outcome; online and on app stores in particular, the purchase is typically a one-time event, meaning that the reputation system becomes fundamental in providing insights on the app quality and making it public for users, in a way to discourage opportunistic behaviours of app developers, even if the transaction is often a unique event (Müller et al., 2011, p. 67). 

Anyway, mobile apps are continuously evolving at a rapid pace with updates being released and replacing older versions quite frequently. As a result of app updates, user satisfaction may change frequently and traditional rating systems appear sometimes unable to capture these dynamics since they are usually resilient to average rating fluctuations when the app has already a substantial number of previous reviews. This system disadvantages new emerging developers, considering that initial low quality of the app and consequently low ratings may cause the impossibility to come back to a good average score later and it could encourage developers to release a new version of a badly reviewed app as a completely new one, trying to reset its reputation (Ruiz et al., 2015, p. 11).

Moreover, another problem of review systems on online marketplaces is ``reputation inflation". In fact, average public feedback scores on these platforms have greatly increased and it looks like customers are reluctant to give a low rating (in particular when the feedback is public) because they fear retaliation or they don't want to harm the seller. In addition, the perceived cost of leaving a negative review increases higher the average rating, pressuring consumers to leave above-average scores (Filippas et al., 2018, p. 30; Horton and Golden, 2015, p. 2). In the long term, this phenomenon makes the reputation system less informative and also creates nonlinear rating expectations for buyers, raising the threshold above which buyers consider the average score a good one.

All in all, the advantages of mobile app stores as platforms and distribution channels are overwhelming, even if some aspects still need to be perfectly tuned. Their introduction contributed to making transaction costs, already decreased by digital economics, even lower and helped in democratizing software distribution.

